
\begin{quote}
El CADi no depende de un vendor único. Cada edición fijará el stack homologado. \ceddie\ cuentas permitidas, residencia de datos y restricciones.
\end{quote}

\subsection*{Categorías}
\begin{enumerate}
\item \textbf{Asistentes y plataformas de agentes:} chats con instrucciones de sistema; entornos con roles, herramientas o flujos multi-paso.
\item \textbf{Edición y colaboración:} documentos colaborativos; Markdown en repositorio para versionado; facilidad de diff humano entre v0 agente y v1 humano.
\item \textbf{Presentaciones:} PowerPoint, Google Slides, Keynote, o generadores desde LaTeX/Beamer/Markdown.
\item \textbf{Gestión de citas:} Zotero, Mendeley, EndNote u homologados --- el humano valida cada referencia.
\item \textbf{Calidad y accesibilidad:} correctores; contrast checkers; checklist del CADi.
\item \textbf{Videoconferencia y aula:} Microsoft Teams u otra plataforma de la edición.
\end{enumerate}

\subsection*{Matriz rápida de uso por sesión}
\begin{tabularx}{\textwidth}{@{}cX@{}}
\toprule
Sesión & Mínimo necesario \\
\midrule
1 & Asistente + documento colaborativo + plantillas \\
2 & Idem + versionado de secciones \\
3--4 & Herramienta de contenido + checklist \\
5--6 & Herramienta de slides + notas \\
7--8 & Todo lo anterior + canal de entrega \\
\bottomrule
\end{tabularx}

\subsection*{Buenas prácticas}
Preferir cuentas institucionales; no pegar datos personales de estudiantes; exportar salidas al agent log; separar experimentación y materiales ``en producción''.

\subsection*{Lo que este CADi no requiere}
Programar frameworks de agentes desde cero; GPU local; comprar software adicional más allá de lo que la institución provea \textit{(salvo decisión explícita de la edición)}.
